\section{2026-10-08 — Ajuste del documento al prototipo construido}

El panel web se describe ahora como lo que se construyó: páginas HTML que genera el propio servicio, una para el histórico del ciudadano y otra para las tendencias de la organización. La organización entra con su clave, que el formulario guarda en una \textit{cookie} \texttt{HttpOnly}, y un administrador da de alta organizaciones y emite o revoca claves con un secreto de configuración. El proveedor de identidad que pedía RF-13 queda como evolución prevista. Las figuras C4 y la de despliegue no se regeneraron: muestran todavía el panel estático y el proveedor de identidad, y sus \textit{captions} lo aclaran.

RNF-15 se da por cumplido con un puntaje SUS medio igual o superior a 68, el valor de referencia de usabilidad aceptable (Sauro, 2011). El canal de supresión del artículo 16 es el correo martinbejarano@gmail.com, informado en la extensión y en el panel, con el plazo de cinco días hábiles; se menciona una sola vez, en la viabilidad legal.

Los casos de prueba 6 y 7 se ejecutaron contra el servicio local sin la credencial del proveedor del LLM y resultaron aprobados con observación: se verificaron clave, cuota, consumo, seudonimización y exportación agregada, pero el análisis llegó parcial y la sección de temas del panel quedó vacía.

\bigskip
